\chapter{Анализ платформы Cosmos}\label{ch:ch2}

\section{Общая архитектура и функциональные возможности Cosmos}\label{sec:ch2/sec1}

Экосистема Cosmos построена как сеть взаимодействующих блокчейнов (Interchain),
ориентированная на создание независимых приложений и их связность через
межсетевой протокол IBC. В отличие от универсальных блокчейнов, Cosmos
поддерживает создание специализированных аппчейнов, каждый из которых сохраняет
суверенитет и может выбирать собственные параметры безопасности и управления.

Исторически развитие блокчейн-платформ пошло от специализированного Биткойна к
универсальному Ethereum. Подход Ethereum дал мощную компонуемость внутри единого
L1, но выявил системные ограничения: ограниченную пропускную способность,
высокие комиссии при нагрузке, рост MEV и сложность апгрейдов, требующих
консенсуса всей сети~\cite{Wood2014,Zhou2019}. Масштабирование через L2-роллапы
снижает комиссии, но фрагментирует ликвидность и усложняет пользовательский опыт.
Cosmos предлагает альтернативную модель: множество суверенных сетей, которые
стандартизированно взаимодействуют через IBC, сохраняя независимость логики и
управления.

Базой архитектуры является Tendermint, обеспечивающий консенсус и сетевой уровень.
Tendermint использует BFT-модель с Proof-of-Stake и позволяет отделить сетевой и
консенсусный уровни от прикладного, чтобы разработчики могли сосредоточиться на
логике приложения. В ядре узла работает CometBFT, который реализует консенсус и
взаимодействует с приложением через интерфейс ABCI.
~\cite{Kwon2014}

ABCI определяет протокол взаимодействия между консенсусным модулем и приложением.
Критически важные методы включают InitChain, BeginBlock, DeliverTx, EndBlock и
Commit, которые обеспечивают корректное обновление состояния. Для проверки
транзакций до включения в блок используется CheckTx. Такой контракт между
уровнями позволяет приложению сохранять независимость, а консенсусу оставаться
стабильным при изменении прикладной логики.

Прикладной уровень реализуется через Cosmos SDK, который предоставляет модульную
архитектуру. Модули отвечают за отдельные функции (например, аутентификация,
стейкинг, токены) и могут подключаться по необходимости, ускоряя разработку
новых сетей. Эта модульность снижает порог входа и позволяет переиспользовать
решения между сетями, что особенно важно в контексте повторного использования
валидаторов.

Для расширения функциональности используются смарт-контракты, например через
CosmWasm. Это дает возможность добавлять бизнес-логику без изменения базового
протокола и упрощает экспериментирование с новыми моделями взаимодействия.

\subsection{Tendermint/CometBFT и мгновенная финальность}\label{sec:ch2/sec1/tendermint}

CometBFT (ранее Tendermint Core) реализует лидер-ориентированный BFT-консенсус для
репликации состояния. На каждой высоте блока протокол согласует следующий блок с
порогом более 2/3 голосующей мощности валидаторов и обеспечивает мгновенную
финальность: как только наблюдается кворум precommit, блок считается окончательно
зафиксированным без вероятностных откатов.

Протокол организован в раунды с последовательностью состояний
NewHeight--Propose--Prevote--Precommit--Commit. Он работает в модели частичной
синхронности: возможны периоды нестабильности сети, но после стабилизации
прогресс гарантируется. Важным элементом является взвешенное голосование:
кворум определяется не числом узлов, а суммарной голосующей мощностью, что
позволяет формально обеспечить безопасность при наличии менее 1/3
злонамеренных валидаторов.

\subsection{IBC и межблокчейновая коммуникация}\label{sec:ch2/sec1/ibc}

IBC задает проверяемую схему взаимодействия между сетями. Каждый блокчейн хранит
on-chain light client другого блокчейна и проверяет входящие данные через
криптографические доказательства (подписанные заголовки и Merkle-проверки).
Релейеры лишь транспортируют сообщения; доверие к ним не требуется. В результате
сохраняется суверенность цепей при наличии стандартизированной совместимости
~\cite{Goes2020}.

\subsection{Детерминизм как базовое требование}\label{sec:ch2/sec1/determinism}

Для BFT-консенсуса критично, чтобы приложение было строго детерминированным:
одинаковый упорядоченный набор транзакций должен давать одинаковый переход
состояния и одинаковый AppHash на всех узлах. В Cosmos это требование фиксируется
на уровне спецификации CometBFT и ложится на разработчика. Нельзя использовать
недетерминированные источники (случайные числа без фиксированного seed, прямое
обращение к системному времени). В среде CosmWasm детерминизм обеспечивается
ограничениями исполнения и контролем API, что снижает риск расхождения состояния.

\section{Преимущества и недостатки платформы}\label{sec:ch2/sec2}

К ключевым преимуществам Cosmos относятся:
\begin{itemize}
  \item модульная архитектура, позволяющая быстро собирать цепи под конкретные
  сценарии использования;
  \item четкое разделение уровней и интерфейс ABCI, упрощающие эволюцию приложений;
  \item поддержка межсетевого взаимодействия через IBC и возможность построения
  экосистемы из суверенных аппчейнов.
\end{itemize}

Среди недостатков и ограничений можно отметить:
\begin{itemize}
  \item значительные операционные затраты при ручной настройке и сопровождении
  валидаторской инфраструктуры;
  \item сложность управления безопасностью в межсетевых сценариях и зависимость
  от выбранной модели защиты;
  \item отсутствие универсального механизма автоматизации отбора валидаторов и
  «матчинга» между валидаторами и потребителями инфраструктуры.
\end{itemize}

Эти особенности задают требования к предлагаемой архитектуре автоматизации,
которая должна учитывать производительность валидаторов, снижать операционные
затраты и повышать управляемость сетей в рамках Cosmos.
